\documentclass[../main.tex]{subfiles}
\begin{document}
%==========================================TIÊU ĐỀ TẬP========================================
\begin{table}[h]
  \begin{adjustwidth}{-1cm}{}
    \begin{tabular}{>{\centering\arraybackslash}p{6cm}|>{\raggedright\arraybackslash}p{12.5cm}}
      \multirow{5}{6cm}
      % Cột bên trái: ÉPISODE và số tập
      {\centering
        {\\[-40pt]
          \flaregothic\fontsize{20pt}{20pt}\selectfont
          \color{eptype}ÉPISODE\color{black}\\
          \burbank\fontsize{72pt}{72pt}\selectfont
          \color{epnum}119
          \\[10pt]
          \hrule
          \vspace{8pt}
          \flaregothic\fontsize{20pt}{20pt}\selectfont
          \color{eptype}SPÉCIAL\color{black}\\
          \burbank\fontsize{72pt}{72pt}\selectfont
          \color{epnum}XXII
          \\[18pt]
          % \flaregothic\fontsize{14pt}{14pt}\selectfont
          % \color{epattr}BIENVENUE, \uline{\color{Banpire} Mel} !
          \flaregothic\fontsize{18pt}{18pt}\selectfont
          \color{gray} COMPILATION\\
          \color{black}
        }
      }
       &
      % Dòng thứ nhất: tên tập tiếng Nhật
      {\fotskip\fontsize{18pt}{18pt}\selectfont \ruby{終}{しゅう}\ruby{焉}{えん}のトリガーになった\ruby{瞬}{しゅん}\ruby{間}{かん}\ 【\topten】
      } \\[6pt]
      % Dòng thứ hai: tên tập tiếng Anh
       & {\cafeteria\fontsize{18pt}{18pt}\selectfont The Beginnings of the End}                                       \\[4pt]
      % Dòng thứ ba: tên tập tiếng Việt
       & {\vnmsans\fontsize{18pt}{18pt}\selectfont Top 10 - Các tập "tận thế" nhất}                                   \\[8pt]
      % Dòng thứ tư: Nhân vật xuất hiện
       & {\sen{\fontsize{14pt}{14pt}\selectfont Track used: }
           \begin{itemize}
             \item {\sen Mutation Crisis (ミューテイション・\ruby{危}{き}\ruby{機}{き}) (Zoftle - Anarchy)}
             \item {\brian Modern Day - Zombotany}
             \item {\ab Angry Birds 2 - Fight and Flight!}
           \end{itemize}
         }
      \\[12pt]
      % Dòng cuối cùng: Thời gian diễn ra của tập (thường sẽ là ngày phát hành)
       & {\continuum\fontsize{16pt}{16pt}\selectfont Ngày phát hành: 22/08/2021}                                      \\[36pt]
    \end{tabular}
  \end{adjustwidth}
\end{table}
%=========================================TRUYỆN CHÍNH========================================
\color{black}

\begin{center}
  (Đây là danh sách top 10 do fan bình chọn cho những tập hỗn loạn nhất của \textit{holo no Graffiti}.)
\end{center}

Mặt trời buổi sớm nhẹ nhàng rọi xuống một thung lũng xanh tươi, nơi tiếng chim hót vang lên giữa khung cảnh núi đồi yên bình. Sora đứng chính giữa khung hình, trong bộ trang phục biểu diễn sáng màu, nụ cười rạng rỡ hiện lên trên khuôn mặt. Gió nhẹ thổi qua, làm mái tóc cô khẽ bay và mang theo làn hương từ rừng thông sau lưng.

Đưa tay nháy mắt tinh nghịch, cô niềm nở nói với khán giả: ``Tóm tắt câu chuyện vừa rồi!''

Nhưng ngay khi lời vừa dứt, bầu trời trong xanh đột ngột chuyển sang đỏ rực như bị đốt cháy. Tiếng rít của lửa và kim loại va chạm vang lên. Cảnh vật phía sau cô biến thành địa ngục trần gian: những tòa nhà cháy rụi, từng mảng tường nổ tung, mặt đất nứt toác và bụi khói cuồn cuộn bốc lên như cột lửa. Những tia chớp đỏ từ dưới lòng đất phóng thẳng lên trời, tạo thành một thứ ánh sáng dữ dội đến lóa mắt.

\begin{figure}[H]
  \centering
  \begin{subfigure}[t]{0.45\textwidth}
    \centering
    \includegraphics[width=8cm]{../../../../Image/Book?/s-22a.png}
  \end{subfigure}
  \quad
  \begin{subfigure}[t]{0.45\textwidth}
    \centering
    \includegraphics[width=8cm]{../../../../Image/Book?/s-22b.png}
  \end{subfigure}
  \caption{Hai thái cực đối lập nhau trong phần mở đầu.}
\end{figure}

Sora vẫn đứng đó, lần này với vẻ mặt nghiêm túc đến mức lạ thường. Không còn nụ cười, chỉ còn ánh mắt băn khoăn và giọng nói chậm rãi, đầy trầm ngâm: ``Bọn mình đã làm gì sai chứ? Tại sao chuyện này lại xảy ra?''

Cô quay mặt nhìn ra phía sau, nơi một tòa nhà nghiêng ngả trước khi sụp xuống như bị kéo vào hư vô. Một ánh đèn sân khấu rọi xuống, tiến theo từng bước chân của Sora trong khi cô nàng tiếp tục trần thuật: ``Mình không thể cứ ngồi đó mà chấp nhận số phận bi đát này được.''

Quay lại phía khán giả, hay đúng hơn là người đang theo dõi đoạn phim, cô nhìn thẳng vào ống kính, ánh mắt dứt khoát: ``Mình phải quay lại những sự kiện trong quá khứ để tìm ra nguyên nhân của mọi chuyện.''

Một khoảng dừng ngắn. Gương mặt cô đột nhiên bừng sáng trở lại như có công tắc bật sẵn, đôi mắt cong cong theo nụ cười và giọng nói vui vẻ như thường lệ: ``Chúng ta hãy cùng xem lại nào!''

Cô nâng một tay lên chỉ thẳng về phía khán giả, giọng cao vút như người dẫn chương trình truyền hình:

\begin{center}
  \textbf{``Top 10 cảnh của \textit{hologra} có khả năng tạo thành ngày tận thế' xin được bắt đầu ngay bây giờ!''}
\end{center}

Và tất cả những lời ấy được cất lên giữa một khung cảnh bốc cháy, với đống đổ nát cháy âm ỉ phía sau, vài mảnh đạn pháo rơi rào rào từ bầu trời và những hình bóng không rõ danh tính chạy hoảng loạn. Sora vẫn giữ nguyên nụ cười tươi rói, một sự vui vẻ đến kỳ lạ giữa cơn đại họa, như thể tất cả chuyện này chỉ là một tiểu phẩm nhẹ nhàng buổi sáng.

\il

\begin{center}
  \uline{(Phần bình luận của Sora sẽ được viết dưới dạng thoại.)}
\end{center}

\pov{Tokino Sora pha lẫn vào ngôi kể thứ ba}

\bxh{10}{Con tàu của Suisei!}{109}{Chín}

[\dots] Cùng lúc đó, Kanata vừa bước vào phòng. Nhìn thấy cảnh tượng trước mắt, cô thoáng sững người, sau đó quay sang hỏi người quản lý đang xoay ghế lại phía mình: ``Tàu\dots\ Suisei? Thế này là thế nào đây, Kuon-senpai?''

Cậu Tổng nhún vai: ``Không biết được, bỗng nhiên em ấy muốn làm con tàu, thế là tự tạo cho mình một con bằng bìa các-tông đây này. Đã chạy vòng quanh phòng như vậy khoảng 15 phút rồi.''

\begin{dialogue}
  \speak{\textit{Sora}} À, cảnh mà Suisei-chan giả vờ làm đoàn tàu đây mà. Trông cô ấy cps vẻ vui đấy, nhưng đó chỉ là bề nổi thôi\dots
\end{dialogue}

Nàng Thiên sứ nhìn tiền bối với ánh mắt không thể tin nổi. Nhưng chưa kịp nói gì, Suisei đã chạy ngay đến bên cô, hai mắt sáng rỡ: ``Em! Kanata, ở đây này!''

Cô nắm lấy tay Thiên sứ và hào hứng hỏi: ``Em có muốn lên chuyến tàu Suisei Cáu kỉnh này không?''

Kanata giật mình: ``Đấy là chị khi cáu kỉnh ư?

Kuon bình thản tiếp lời: ``Có thể em không tin, nhưng đúng là như vậy.''

Cô nàng thế hệ thứ Tư vẫn chưa hết bàng hoàng, liền hỏi tiếp ``tàu hỏa'' Suisei: ``Không phải chị thường hay đi hát sao?''

\begin{center}
  \uline{(Có một cái đồng hồ đếm ngược ở góc dưới màn hình. Đồng hồ chỉ về không khi cảnh ở dưới xảy ra.)}
\end{center}

Nhưng ngay lập tức, cô nhận lại một ánh mắt sắc lạnh từ nàng Ca sĩ. Cô nàng trừng mắt, giọng đơn điệu và đầy vẻ triết lý: ``Không có niềm vui nào là kéo dài mãi chứ?''

\img{9-1b}

``Mặt chị nhìn sợ quá!!'' cô nàng hậu bối run rẩy, mặt tái mét.

Còn Kuon thì bật cười: ``Đây mới chính là cáu kỉnh này. Mà Suisei-chan, đừng làm em ấy sợ.''

\begin{dialogue}
  \speak{\textit{Sora}} Đúng như các bạn thấy. Cơn tức giận của cậu ấy quả thực có khả năng hủy diệt cả thế giới.
\end{dialogue}

Vẻ mặt của nàng tóc xanh nhanh chóng vui tươi trở lại. Cô vỗ nhẹ lên tấm bìa các-tông và nói với hậu bối của mình: ``Nói thế thôi, em cũng lên thử đi. Cùng lên tàu với Sui-chan nào!''

Kanata do dự nhìn sang cậu Tổng, ánh mắt như muốn hỏi ``Anh có chắc là bình thường không đấy?''

Đáp lại, cậu chỉ cười khì và ra dấu hiệu: ``Cứ lên tàu đi. Anh đang bận việc, chắc sẽ không lên đâu.''

Nàng Thiên sứ miễn cưỡng trèo lên con ``tàu hỏa'', trong khi chủ tàu hân hoan tuyên bố: ``Con tàu Suisei đã sẵn sàng!''

\begin{dialogue}
  \speak{\textit{Sora}} Mọi người đã hiểu rồi chứ? Sự kiện ở tập đó đã diễn ra như vậy đấy!
\end{dialogue}

``\textit{Và chúng ta có một con tàu Suisei, với `hai mẹ con' ở trên đó,}'' Kuon cười thầm bình luận trong lòng.

``Cái quái gì đây\dots'' Kanata nhăn mặt, bày tỏ sự khó hiểu với thứ bòng bong mà cô đang dây vào [\dots].

\bxh{9}{Cừu uống cà phê!}{89}{Bảy}

[\dots] Cô vùng lên, đập tay xuống bàn, thốt lên một câu đầy khẳng định: ``Em sẽ trở thành một cô Cừu có thể uống cà phê!''

Kuon phì cười, chống cằm nhìn cô đầy thích thú: ``Vừa nãy uống còn phun nó vào mặt Subaru-chan thế kia mà còn đòi.'' Cậu quay sang Subaru, nhếch mép cười.

``Mà nói đến Subaru-chan thì\dots''

Ngay giây tiếp theo, một luồng sát khí từ Subaru tỏa ra phía sau lưng Watame. Cô nàng Vịt đặt tay lên vai nàng Cừu, giọng nói trầm xuống một cách nguy hiểm. Cô siết nhẹ vai nàng Cừu, cất lên một giọng đáng sợ ``\dots Xin lỗi chị đi đã.''

\begin{dialogue}
  \speak{\textit{Sora}} Đây là một cảnh mà Watame-chan sau khi phun cà phê đen vào mặt của Subaru-chan\dots
\end{dialogue}

Watame cứng đờ người.

Trong khi đó, Choco giả vờ không thấy cảnh tượng đó, tiếp tục hỏi với giọng điệu đầy trêu chọc: ``Rồi, rồi. Em có làm được không đấy, Watame-sama?''

Watame sà tới với vẻ mặt cực kỳ nghiêm túc: ``Em làm được! Em biết mà!''

\begin{center}
  \uline{(Có một cái đồng hồ đếm ngược ở góc dưới màn hình. Đồng hồ chỉ về không khi cảnh ở dưới xảy ra.)}
\end{center}


Vừa dứt lời, cô lập tức bị Subaru bẻ người nhẹ ra phía sau. Nàng Cừu giãy nảy, cố cầm hai cánh tay đang ghì chặt cô: ``Uwawawwa\dots\ Tha cho emmmmmm!!''

Choco nhấp một ngụm cà phê rồi nói tỉnh bơ: ``Vậy em nên bắt đầu bằng việc hiểu rõ xem mình đang ở đâu đi.''

\img{7-7f}

Nàng Cừu vẫn đang vẫy vùng để thoát khỏi Subaru, người vẫn siết chặt như một chiếc khóa gọng kìm.

\begin{dialogue}
  \speak{\textit{Sora}} Mình nghe nói là có một đấu vật siêu hạng nào đó bắt đầu sử dụng đòn này trong trận đấu của người đó rồi đấy! Nhưng mà các bạn à, đừng làm thử tại nhà nha.
\end{dialogue}

Cuối cùng, sau một hồi vật lộn, Watame thoát ra được, đứng thẳng người, đập tay xuống bàn một cái *bốp!*: ``Chúng ta cùng bắt đầu cuộc họp chiến thuật nào!''

Không khí trong phòng đột nhiên biến thành một cuộc họp quân sự.

Watame quay sang nhìn Subaru, người đang lấy khăn lau cà phê vẫn còn dính trên mặt. Cô chỉ tay thẳng về phía nàng Vịt: ``Đi thôi nào, Subaru-senpai!''

Subaru giật mình, tròn mắt nhìn cô nàng Cừu đầy cảnh giác: ``Ể\dots\ Chị không muốn đi-!?'' Nhưng trước khi cô kịp phản ứng, thì hậu bối đã kéo cô đi ra khỏi văn phòng [\dots].

\bxh{8}{Chiếc đồng hồ nổi giận!}{100}{Tám}

[\dots] ``PolPol tập trung vào nha\~{}, OK?'' - nàng Cáo sa mạc hỏi với cái miệng to đến mức bất thường, rồi đứng lên, bấm nút phóng.

``Ta biết quá rõ thế này rồi, nhưng cảm ơn!'' - Polka hét lên khi bị bắn vèo đi như một viên đạn sống.

\begin{dialogue}
  \speak{\textit{Sora}} Đây là cảnh mà bọn họ (chỉ bộ ba Aqua-Marine-Kuon) bị đuổi bắt bởi nhóm băng đảng trong rừng.
\end{dialogue}

``\dots Mẹ kiếp.'' Đó là những lời nói cuối cùng của Kuon trước khi ``thủ lĩnh'' của băng đảng  đáp xuống chiếc xe Mio, tạo ra một vụ nổ nhẹ khiến mọi người bị hất văng xuống đất như bắp rang bị thổi tung trong lò vi sóng.

Miofast hét lên thảm thiết, mất lái và đâm sầm vào chiếc đồng hồ một cách định mệnh, rồi phát nổ như một bộ phim hành động kinh phí thấp.

\begin{center}
  \uline{(Có một cái đồng hồ đếm ngược ở góc dưới màn hình.)}
\end{center}

Tất cả mọi người, trừ Marine và Kuon ra, đều đã ngất xỉu hoàn toàn. Hai người còn tỉnh táo đứng dậy, quần áo bám đầy bụi, trong khi chiếc xe hình Sói đen cùng chiếc đồng hồ biết chạy kia đã cùng nhau ``tử nạn''.

``Cái quái gì vừa xảy ra vậy\dots? Ái da\dots'' Kuon vừa xoa đầu, vừa lồm cồm ngồi dậy.

Marine đưa tay về phía luồng khói, như một nhân vật chính trong bộ phim mất đi người bạn thân trong chiến tranh: ``Đồng hồ ơi\dots\ Không\dots''

Cậu Tổng đứng dậy phủi bụi, gắt gỏng: ``Nó bị giết chết rồi. Tuyệt cmn vời.''

Mặc kệ những lời bực bội của cậu, nàng Thuyền trưởng vẫn rơm rớm nước mắt, nghẹn ngào trước chiếc đồng hồ cầm tay mà chính cô đã dọa nạt trước đó: ``Ta sẽ nhớ đến ngươi, đồng hồ ơi\dots''

``Anh bắt đầu phát điên với nơi này rồi đấy\dots'' cậu thở dài, thi thoảng ho vài tiếng.

Khói bụi trước mặt họ vẫn bốc lên nghi ngút như thể đang che giấu một thứ gì đó\dots

\begin{center}
  \uline{(Đồng hồ chỉ về không khi cảnh ở dưới xảy ra.)}
\end{center}

``Hửm?'' / ``Cái quái\dots''

Và đó là lúc mà linh hồn tức giận của cái đồng hồ vùng lên. Theo lời kể của Miko.

Từ trong màn khói, một bóng dáng kỳ dị dần hiện ra. Đó là một người có thân hình giống A-chan, nhưng thay vì một khuôn mặt bình thường, gương mặt của cô là một cái đồng hồ. Nó lao tới với tư thế như đang đấm bay hai người kia, miệng lặp đi lặp lại câu: ``Trời sáng rồi mấy đứa ơi! Dậy đi nào!''

\img{s-18o}

\begin{dialogue}
  \speak{\textit{Sora}} Giá như linh hồn giận dữ ấy không chiếm lấy thân xác của chiếc đồng hồ đó\dots\ Nhưng rất may cho họ đó chỉ là một giấc mơ mà thôi.
\end{dialogue}

Marine và Kuon tròn mắt, chưa kịp phản ứng, cái mặt đồng hồ đã tung cước húc cùi chỏ thẳng vào đầu cả hai, khiến họ bay ngược ra sau như cảnh tua chậm trong phim hành động.

Trong giây phút mất ý thức, họ chỉ có thể nghe thấy một giọng nói vang vọng của bà thần Okayu: ``Tạm biệt, nyan\~{}''

Và rồi\dots\ màn đêm buông xuống [\dots].

\bxh{7}{Tòa nhà Kuon!}{28}{Hai}

[\dots] Ngay khi cậu vừa dứt câu, ``Người dưới đất'' lại bất ngờ trồi lên, lần này tóm chặt lấy chân Kuon. Hắn ta nói với giọng lạc hẳn đi: ``Cậu cũng sẽ phải làm phân bón nữa.''

``\textbf{Khoan đã! Cái gì cơ!?}'' cậu hét lớn, nhưng rồi hắn ta đã lôi cậu xuống đất trước sự bàng hoàng của cả Korone và Okayu.

\begin{dialogue}
  \speak{Korone} Kuon-senpai!?
  \speak{Okayu} Kuon-senpai à!?
  \speak{\textit{Sora}} Đây là cảnh mà Kuon-senpai bắt đầu khởi nghiệp này!
\end{dialogue}

\begin{center}
  \uline{(Có một cái đồng hồ đếm ngược ở góc dưới màn hình. Đồng hồ chỉ về không khi cảnh ở dưới xảy ra.)}
\end{center}

Chưa đầy một giây sau, khuôn mặt của Kuon bất ngờ hiện lên trên bề mặt cái mầm cây. Và rồi, từ cái mầm nhỏ bé đó, một tòa nhà khổng lồ mọc lên, vươn cao vút trời xanh.

\begin{figure}[H]
  \centering
  \begin{subfigure}[b]{0.45\textwidth}
    \centering
    \includegraphics[width=8cm]{../../../../Image/Book/2-12g1.png}
  \end{subfigure}
  \quad
  \begin{subfigure}[b]{0.45\textwidth}
    \centering
    \includegraphics[width=8cm]{../../../../Image/Book/2-12g2.png}
  \end{subfigure}
  \quad
  \begin{subfigure}[b]{0.45\textwidth}
    \centering
    \includegraphics[width=8cm]{../../../../Image/Book/2-12h.png}
  \end{subfigure}
  \caption{Quá trình từ mầm cây mọc lên tòa nhà.}
\end{figure}

Sự xuất hiện đột ngột của tòa nhà cũng vô tình đưa Korone lên đỉnh. Bạn của cô đứng từ dưới, mắt cô sáng quắc ngước nhìn lên, dang rộng hai tay và reo lên đầy hào hứng: ``Công ty Kuon đã được ra đời! Koro-san đứng trên cao kìa\dots!!''

Korone đứng trên đỉnh tòa nhà, vẫy tay vui vẻ: ``Wai\~{}! Trên này cảnh hùng vĩ quá\~!{}''

\begin{dialogue}
  \speak{\textit{Sora}} Công ty Kuon hẳn sẽ làm ăn phát đạt lắm đây.
\end{dialogue}

Nàng Khuyển, không quên bản chất năng động của mình, nhảy tưng tưng trên đỉnh tòa nhà mới xây. Trong khi đó, nàng Miêu cũng dang tay múa may dưới chân tòa nhà. Cặp đôi lại tiếp tục màn nhảy múa tưng bừng của mình, mặc kệ mọi sự kỳ lạ vừa xảy ra [\dots].

\bxh{6}{Cơn tức giận của Shion!}{34}{Ba}

Cả cậu và nàng Hiệp sĩ kinh ngạc, trong khi nàng Thỏ chỉ ngơ ngác nhìn cô đầy khó hiểu. còn nàng Chó nâu thì reo lên: ``Đó! Em đã bảo rồi mà!''

\begin{dialogue}
  \speak{Noel} Tại sao chị lại bị kẹt trong tường!? Chuyện này là sao hả, Shion-senpai?
  \speak{Kuon} Với cả, mấy cuốn catalog kia là sao?
  \speak{Korone} Trông dáng bà ngố tàu thật á\~{}
  \speak{\textit{Sora}} Chắc cái này mình không cần giải thích thêm nữa nhỉ.
\end{dialogue}

\begin{center}
  \uline{(Có một cái đồng hồ đếm ngược ở góc dưới màn hình.)}
\end{center}

Nàng Phù thủy nhìn cả nhóm, khuôn mặt vừa xấu hổ vừa cố tỏ ra bình thường: ``À thì\dots\ chị đang luyện dịch chuyển tức thời\dots\ nhưng mà hơi trượt tay nên là\dots''

Cậu khoan tay, nhíu mày nghiêm nghị: ``\dots kẹt trong tường như này hả? Thế còn mấy cuốn catalog thì sao? Anh không nghĩ chúng tự chui vào tường đâu.''

Shion cười gượng: ``Ừm\dots\ thì\dots\ em thực hành với chúng đó. Ai mà ngờ nó lại thành công đến mức\dots\ vượt mong đợi như này!''

Kuon thở dài, lắc đầu ngao ngán: ``Thành công tới mức kẹt tường thế này\dots\ Em làm anh hết nói nổi rồi đấy, Shion-chan ạ.'' Cậu bước tới, định kéo nàng phù thủy ra khỏi bức tường, nhưng rồi hỏi: ``Mà này, thế quái nào em có thể kẹt được như thế? Cả đầu, cả tay chân đều bị ép sát thế này á!?''

``Làm sao mà em biết được!'' nàng Phù thủy bật lại, ``Em còn không hiểu sao nó thành ra thế này nữa!''

Cả nhóm chỉ biết nhìn nhau, nửa buồn cười, nửa bối rối trước sự vụ không thể nào kỳ lạ hơn này.

\begin{center}
  \uline{(Đồng hồ chỉ về không khi cảnh ở dưới xảy ra.)}
\end{center}

Nàng Thỏ nhìn cảnh tượng oái oăm của đàn chị mình, không giấu được nụ cười mỉa mai: ``Đúng là ngu thật mà, peko.''

Đến cậu Tổng cũng phải đồng tình với cô: ``Peko-chan nói đúng đấy. Không cần phải bàn cãi.''

Shion ngay lập tức nhăn mặt, giọng pha chút bức xúc: ``Ý anh là sao hả, Kuon-senpai!? Mà còn em nữa, Pekora-chan, không phải em cũng có cà rốt kẹt trong tóc đấy thôi?''

\begin{dialogue}
  \speak{\textit{Sora}} Có lẽ nếu Shion-chan không bị mọi người dồn dập trêu chọc thế này, em ấy sẽ không phát hoảng đến như vậy\dots
\end{dialogue}

\img{3-5c}

``Hai cái đó khác nhau hoàn toàn về bản chất đó em ạ,'' cậu thở dài, ``một cái là do tai nạn ngu ngốc, một cái là do người ta tự nguyện.''

Nàng Thỏ hừ một tiếng, cảm thấy bị ``đụng chạm lòng tự ái'', liền cất lời khoe khoang: ``Kuon-senpai nói chí phải, peko! Không giống như chị Shion-senpai hậu đậu đâu, mấy củ cà rốt của em có thể gỡ ra bất cứ lúc nào theo ý muốn, peko!''

``Thế á? Vậy em thử làm đi cho bọn anh xem,'' cậu nhướng mày, tỏ vẻ không tin.

Pekora hất mặt, tay khéo léo gỡ củ cà rốt từ trong lọn tóc ra để chứng minh [\dots].

\bxh{5}{Ayame bất ngờ!}{81}{Sáu}

\begin{dialogue}
  \speak{\textit{Sora}} Cái này khỏi phải chờ!
\end{dialogue}

\begin{center}
  \uline{(Chiếc đồng hồ đếm ngược về không ngay lập tức.)}
\end{center}

Cảnh quay bắt đầu với một chiếc ti-vi tĩnh lặng, không có lấy một bóng người.

Đột nhiên, từ dưới khung hình, nàng Quỷ sừng bật lên, mặt tươi rói, reo:

``\textbf{Yo-dayo!}''

\gif{ayame1s}{1}{32}

\begin{dialogue}
  \speak{\textit{Sora}} Hẳn là sẽ có một làn sóng tweet ập vào đây.
\end{dialogue}

Cảnh chuyển nhanh sang một góc khác trong phòng. Ayame vừa đi vừa nói, giọng đầy hào hứng: ``Dạo gần đây trời cũng bắt đầu lạnh rồi nhỉ?''

Đằng sau ống kính, Kuon thầm nghĩ: ``\textit{Thì bây giờ chuẩn bị vào đông rồi mà.}''

Bất thình lình, cô nàng dí sát mặt vào máy quay, đôi mắt ánh lên vẻ tinh nghịch: ``\textbf{Đúng chứ?}''

Cậu giật mình, vội vàng giữ chắc máy quay để không bị rung.

Cảnh tiếp tục chuyển đến một góc khác, nơi nàng Quỷ sừng đang nằm sấp trên sàn, khẽ lắc lư đầu qua lại. Giọng cô vẫn tràn đầy năng lượng: ``Mọi người cũng đồng ý với mình chứ?'' [\dots]

\bxh{4}{Một thùng-LCL-để thở!}{91}{Bảy}

\pov{Tokino Sora xen kẽ với Hikawa Kuon}

\begin{dialogue}
  \speak{\textit{Sora}} Cái này cũng khỏi phải chờ! Cẩn thận nha!
\end{dialogue}

\begin{center}
  \uline{(Chiếc đồng hồ đếm ngược về không ngay lập tức.)}
\end{center}

[\dots] Bên trong căn phòng ấm cúng, lò sưởi đang cháy tí tách, tỏa ra ánh sáng cam dịu nhẹ. Dẫu vậy, bầu không khí lại chẳng hề yên bình như vẻ ngoài của nó.

``Rừm! Rặc! Rặc! Rặc! Rặc! Rặc! Rặc!'' tiếng của Lamy-chan vang lên, phá tan cả sự im lặng vốn có.

\gif{7-9i}{1}{82}

Tôi thở dài, chống tay lên trán: ``Dừng lại đi em ơi.''

Tôi nhìn nàng Yêu tinh tuyết đang vừa rung người, vừa ngước mắt lên trần nhà, miệng không ngừng phát ra những âm thanh khó hiểu. ``Em kêu nãy giờ được nửa tiếng rồi đấy. Dừng lại đi.''

Nàng ta chẳng có vẻ gì là nghe thấy tôi nói, vẫn tiếp tục tạo những âm thanh kỳ quặc kia.

Mio-chan và Flare-chan, hai nạn nhân vô tội của tình huống này, đã lặng lẽ trốn ra sau kệ tủ, chỉ ló mỗi cái đầu ra để quan sát từ xa.

\img{7-9b}

Nàng Dị giới cau mày: ``Lamy-chan đang làm cái quái gì thế?''

Nàng Sói đen liếc đồng đội, giọng thản nhiên như thể đã quen với cảnh tượng này: ``Chắc là đang bắt chước tiếng máy cắt cỏ vì đang chán ấy mà.''

Tôi nhíu mày: ``Lúc buồn chán bộ em ấy toàn làm thế suốt à\dots?''

Flare-chan thở dài: ``Có đầy thứ mà em ấy có thể bắt chước được\dots''

Tôi và Mio-chan nhìn nhau, cùng đồng thanh chốt hạ: ``\dots mà lại bắt chước cái thứ của nợ đó. Lạ ghê.''

\begin{dialogue}
  \speak{\textit{Sora}} Chính điều này đã làm cho máy cắt cỏ bán đắt hàng như tôm tươi, dẫn tới hủy hoại hết cây xanh trên Trái Đất này\dots
\end{dialogue}

Sau một hồi phát ra những âm thanh vô nghĩa đến mức tôi bắt đầu tự hỏi liệu em ấy có bị ``lỗi hệ thống'' không, nàng Yêu tinh tuyết cuối cùng cũng chịu dừng lại. Cô thở phào, lấy ra một chai rượu từ đâu đó rồi tự lẩm bẩm: ``Được rồi, nghỉ một chút nào [\dots].''

{\color{bronze}\bxh{3}{Tiếng kẻng!}{74}{Sáu}}

\pov{Tokino Sora pha lẫn vào ngôi kể thứ ba}

[\dots] Chủ nhân của chiếc kẻng - nàng Dị giới - xuất hiện. Cô tiến đến gần Matsuri, người đang cầm chiếc kẻng trong tay, và hỏi với giọng nhẹ nhàng: ``Em có chơi được không, Matsuri-chan?''

Cô nàng Lễ hội không trả lời mà chỉ nở một nụ cười thật tươi, mắt nhắm lại, như thể là một bức tượng phật trong trạng thái hỉ lạc vĩnh cửu.

Kuon nhìn khuôn mặt hỉ đó, thở dài bất lực: ``Em ấy đã như thế khoảng 15 phút rồi''

Sói đen, người đã tận mắt chứng kiến toàn bộ quá trình, cúi gằm mặt xuống tỏ vẻ mệt mỏi: ``Tâm của chị ấy trong sáng đến mức mỉm cười 24/7 luôn rồi\dots''

Aki tròn mắt ngạc nhiên, hốt hoảng: ``Ehh!? Sợ  thế!''

\begin{dialogue}
  \speak{\textit{Sora}} Matsuri-chan đã mượn được chiếc kẻng ma thuật từ Aki-chan. Tùy thuộc vào độ thuàn khiết của tâm người đánh mà tiếng của nó sẽ thay đổi!
\end{dialogue}

``Ít nhất thì cái tiếng kẻng chắc phải nghe hay hơn rồi nhỉ?'' cậu chống hông, nhìn vào chiếc kẻng.

Mio đứng dậy, đặt tay lên vai Matsuri, rồi hướng ánh mắt về chiếc kẻng trong tay cô nàng. ``Vậy em chơi thử đi!''

Aki hào hứng vỗ tay. ``Được rồi! Một, hai\dots''

Mio, Aki và Kuon bắt đầu đếm nhịp: ``Hai, ba!''

\begin{center}
  \uline{(Có một cái đồng hồ đếm ngược ở góc dưới màn hình. Đồng hồ chỉ về không khi cảnh ở dưới xảy ra.)}
\end{center}

Nàng Lễ hội nâng chiếc gậy lên, rồi gõ mạnh một cái vào kẻng.

*KENGGGGGGGGGGGGGGGGGGGGGGGGGGGGGGGGGGG!!*

\gif{6-5j}{1}{152}

Lập tức, một thứ âm thanh vang vọng cả văn phòng, lan ra khỏi tòa nhà, len lỏi đến thành phố, và rồi cả Trái Đất.

Âm thanh trong trẻo và vang vọng đến mức người viết câu chuyện này ở Hà Nội còn có thể nghe thấy (!?).

\ct

Khoảnh khắc đó, Aki và Mio cũng đột nhiên mỉm cười y như Matsuri - một nụ cười vô thức nhưng lại mang theo sự thanh thản lạ thường.

Trong khi đó, Rushia và Kuon đứng chôn chân tại chỗ, trố mắt nhìn ba người bọn họ với sự ngỡ ngàng và khó hiểu.

\begin{dialogue}
  \speak{\textit{Sora}} Nhóm `Chị em Hút bụi Ion âm' đã trở thành hiện tượng ở thời điểm đó!
\end{dialogue}

Aki mỉm cười rạng rỡ, tuyên bố với một sự phấn khích đến khó tin: ``Và với tiếng kẻng đó, bọn chị đã lập nhóm `Chị em Hút bụi Ion âm' hoàn chỉnh [\dots]!''

{\color{silver}\bxh{2}{Sức mạnh mùa xuân!}{106}{Tám}}

\pov{Tokino Sora xen kẽ với Hikawa Kuon}

[\dots] Cả hai chúng tôi đồng loạt quay về hướng phát ra âm thanh ấy. Một bóng người dần hiện lên từ hư vô.

``Đây là vùng đất mùa xuân, và ta là nữ thần mùa xuân đây.''

Giọng nói này nghe quen quá\dots\ Tôi và Nene-chi chỉ tay thẳng vào người đó, không hẹn mà cùng đồng thanh: ``À! Miko-senpai/-chan kìa.''

Cái bóng lập tức bật dậy, hét lên đầy tức tối: ``\textbf{Này! Đừng có nói toẹt ra như vậy chứ, ne!}''

Tôi nhún vai đáp lại: ``Thì giọng bẹt của em nghe dễ đoán mà. [\dots] Ngay sau đó, tôi nghe thấy một tiếng đập bàn cực mạnh, kèm theo tiếng rít khó chịu của míc-rô vọng lại. Có vẻ như em ấy đã bị tôi bắt bài.

Tôi vội vàng bịt tai lại, còn đồng đội của tôi thì bật cười khúc khích. Miko-chi đúng là dễ bị trêu thật.

\begin{dialogue}
  \speak{\textit{Sora}} Trong công cuộc tìm kiếm ``mMùa xuân'' của mình, Nene-chan vô tình ở trong một không gian kỳ lạ bí ẩn\dots
\end{dialogue}

Sau một hồi trấn tĩnh, Miko- à không, ``nữ thần mùa xuân'' nghiêm giọng tuyên bố, vẫn với giọng bẹt và âm ``ne'' ấy: ``Vì nhà ngươi có niềm tin vào mùa xuân, ta sẽ ban cho ngươi sức mạnh của mùa xuân!''

Nene-chi chớp mắt đầy ngạc nhiên, còn tôi thì cau mày: ``Khoan, sức mạnh mùa xuân là cái gì-''

Chưa kịp nói hết câu, cả người cô nàng Hải cẩu đột nhiên tỏa sáng với một thứ ánh sáng màu hồng rực rỡ. Em ấy ngạc nhiên nhìn hai tay mình phát sáng, còn tôi chỉ có thể đứng đó với một dấu chấm hỏi to đùng trên đầu rẳng thế tại sao lại là em ấy.

``Nữ thần'' kia để lại lời nhắn cuối cùng trước khi biến mất: ``Hãy cho họ biết thế nào là lễ độ.''

Nene-chan siết tay lại đầy quyết tâm, mắt sáng lên: ``Em sẽ cố gắng ạ, Miko-senpai!''

\begin{center}
  \uline{(Có một cái đồng hồ đếm ngược ở góc dưới màn hình.)}
\end{center}

Nhưng ngay lập tức, giọng nói của nữ thần vang lên đầy tuyệt vọng: ``\textbf{Đã bảo rồi! Ta là nữ thần mùa xuân cơ mà, ne!}''

Tôi bật cười, Nene cũng cười theo. Chúng tôi cười một tràng dài vì cái cách Miko-chi phản ứng khi bị chính chúng tôi bắt tại trận quá dễ thương.

Rồi, một ánh sáng màu hồng lại bao trùm lấy cả hai chúng tôi, những dải hoa anh đào rơi xuống, quay một vòng bao quanh chúng tôi.

\ct

\begin{center}
  \uline{(Đồng hồ chỉ về không khi cảnh ở dưới xảy ra.)}
\end{center}

Khi tôi nhận thức lại thì\dots\ chúng tôi đã trở về văn phòng.

Tôi được đưa trở về trước, vừa kịp lấy lại ý thức thì đã nghe thấy hai giọng nói đồng thanh vang lên: ``A, Kuon-senpai về rồi kìa!''

Flare-chan và Kana-chan cùng lúc thốt lên, cùng tỏ ra ngạc nhiên khi thấy tôi đột nhiên xuất hiện giữa phòng.

``Anh vừa đi đâu thế?'' nàng Dị giới hỏi.

Tôi chỉ đáp gọn: ``Đuổi theo Nene-chan. Tìm thấy em ấy rồi.''

Nàng Thiên sứ định hỏi tôi em ấy ở đâu, nhưng chưa kịp cất lời, một luồng ánh sáng hồng chói lóa lại bùng lên ngay phía sau tôi, và một giọng nói đầy kiêu hãnh vang lên: ``Hãy nhìn Super-Nene ta đây, giờ với sức mạnh mùa xuân mới!''

Nene xuất hiện với một tư thế cực ngầu - ít nhất là trong suy nghĩ của em ấy: Nene-chi từ đâu xuất hiện sau luồng ánh sáng hồng hào kia, khẽ lơ lửng trên không rồi nhẹ nhàng chạm đất. Tôi nhíu mày nhìn cô nàng, trong đầu đầy dấu chấm hỏi.

``Rồi, vậy sức mạnh mới đó là gì?''

Nhưng chưa kịp nghe câu trả lời, tôi đã thấy Nene đột nhiên tạo dáng, giơ hai tay lên và bắn ra một thứ tia sáng màu hồng thẳng vào đầu Thiên sứ.

\img{8-10f}

Tôi vội vàng chuẩn bị tinh thần để đỡ một trận la hét vì đau đớn\dots\ nhưng trái ngược với dự đoán, Kana-chan lại tròn mắt, tỏ ra bất ngờ: ``A! Chị hết đau đầu rồi này!''

Tôi chớp mắt ngạc nhiên. Hả? C-Cái gì?

Nene-chi cũng không kém phần kinh ngạc, em ấy nhìn bàn tay mình rồi nhìn nàng Thiên sứ với vẻ không tin nổi.

Tôi sửng sốt tới mức toát mồ hôi hột: ``V-Vậy cũng được nữa?''

Ngay lúc đó, một cơn gió nhẹ bất ngờ lướt qua căn phòng. Nhưng cơn gió này\dots\ lại đến từ một thực thể có hình dạng của một cánh hoa đào, lại còn được vẽ bằng giấy!? Nó trôi lơ lửng một cách đầy ma mị.

\begin{dialogue}
  \speak{\textit{Sora}} Thế là khái niệm ``bốn mùa luôn thay đổi'' kết thúc tại đây.
\end{dialogue}

Flare liền hít một hơi sâu, đôi mắt ánh lên vẻ xúc động: ``Cơn gió xuân đầu tiên kìa!''

Tôi khoanh tay lại, cau mày nhìn cánh hoa đào giấy màu hồng kia: ``Sao anh cảm thấy có gì đó sai sai nhỉ?''

Nene-chi không trả lời, chỉ búng tay một cái, khiến cho cơn gió kỳ lạ ấy bay đi mất. Tôi bắt đầu cảm thấy tình huống này ngày càng kỳ quặc.

Thế rồi, tôi chưa kịp suy nghĩ thêm về tình huống ky quặc đó thì đã nghe thấy hai người con gái kia đồng loạt lên tiếng bằng một giọng điệu đầy ngưỡng mộ:

\begin{dialogue}
  \speak{Kanata} Nene-chan, cảm ơn em. Em là vị cứu tinh đấy.
  \speak{Flare} Một nữ thần! Là một nữ thần mùa xuân đó!
\end{dialogue}

Tôi chỉ biết lắc đầu nhìn tình huống khó hiểu này, buột miệng: ``\dots Quá sai luôn.''

{\color{gold}\bxh{1}{Mảnh thiên thạch!}{82}{Sáu}}

\pov{Tokino Sora pha lẫn vào ngôi kể thứ ba}

[\dots] Haachamasaurus, với toàn bộ sức mạnh của mình, ném mạnh chiếc đầu (Fubuki) giả lên trời với tốc độ siêu thanh như một quả bóng chày.

``\textbf{HAACHAMA-CHAMA!!}''

*VÚT!*

\begin{dialogue}
  \speak{\textit{Sora}} Cảnh mà Haato-chan thoát khỏi hang ổ của Subaru-chan và quay trở lại bầy đàn với một tâm trạng bực dọc.
\end{dialogue}

Chiếc đầu bay thẳng về phía chiếc PC của Rushiasaurus - sau đây gọi là Rushia - kẻ vốn chẳng quan tâm đến trận đấu mà vẫn đang tập trung vào trò chơi của mình.

*RẦM!* *BÙM*

Cú va chạm cực mạnh. Chiếc thùng máy tính phát nổ, tạo ra một vụ chấn động kinh hoàng, hất văng tất cả những con khủng long xung quanh ra xa.

``Có vẻ như đến cả những người không liên quan cũng bị vạ lây.'' Kuon tiếp tục bình luận.

\begin{center}
  \uline{(Có một cái đồng hồ đếm ngược ở góc dưới màn hình.)}
\end{center}

Rushiasaurus bàng hoàng nhìn về phía PC của mình. Màn hình hiển thị dòng thông báo lỗi, ứng dụng bị văng ra, toàn màn hình chỉ còn lại một màu xanh lè với hàng đống dòng lỗi và một thông báo như đang trêu người cô. Cô đồng hoảng loạn cực độ.

\begin{center}
  \uline{(Đồng hồ chỉ về không khi cảnh ở dưới xảy ra.)}
\end{center}

*RẦM!*

Cô đập mạnh xuống bàn. ``\textbf{Tôi đã làm gì nên tội chứ!? Tại sao lại là tôi!?}''

\img{6-13s}

{\Large \textbf{*RẦM!*}}

Cú đập thứ hai mạnh đến mức gây ra một vụ nổ kinh hoàng, thổi bay mọi thứ xung quanh. Một vùng Nhật Bản chìm trong ánh sáng trắng chói lòa.

Khi ánh sáng tan đi, thế giới đã thay đổi. Mọi thứ đều trở nên khác lạ, nhưng cũng không hẳn là hoàn toàn xa lạ.

``Và thế là, kỷ Phấn trắng đã đến hồi chấm dứt.'' Fubuki tiếp lời cậu con trai với một giọng điệu đầy nuối tiếc.

``\hll\ đã tồn tại kể từ sự bùng nổ của kỳ Cambri, nhưng giờ, chúng đã rơi vào trạng thái ngủ đông.''

Ngoài không gian, một góc của Trái Đất lóe sáng như thể vừa xảy ra một vụ nổ có sức hủy diệt ngang với bom nguyên tử.

\img{6-13t}

``\hll\ chỉ nổi lên tiếp tục và phát triển từ thời kỳ Đại Tân sinh,'' Kuon tiếp lời cô, kết thúc bình luận của mình.

Bề mặt Trái Đất giờ đây chỉ còn lại một bờ cát trắng, lác đác vài cây dừa. Những người sống sót xuất hiện.

Sora đang ngủ.

Rushia ngồi bệt xuống đất, tay vẫn còn run rẩy.

Haato thì bị cắm đầu xuống cát, hai chân giãy giãy.

Aki, Miko và Choco thì từ biển bước lên bờ, toàn thân ướt sũng.

Ở xa xa trên bầu trời, một hình bóng quen thuộc vẫn đang lởn vởn: linh hồn của Subarusaurus, ánh mắt đầy oán hận, nhìn chằm chằm xuống mặt đất như thể đang hăm dọa những kẻ đã tiêu diệt mình.

\img{6-13u}

\begin{center}
  -Kết thúc.-
\end{center}

\il

Cảnh quay trở lại với khung cảnh tận thế ban đầu.

Gió thét gào như tiếng hú của một thế giới đang hấp hối. Những tia lửa từ mặt đất vẫn không ngừng bắn lên, hòa vào bầu trời đỏ quạch như máu đang sôi. Trong không gian đổ nát ấy, Sora đứng run rẩy giữa những mảnh gạch vỡ và khói bụi mù mịt.

Nàng Thần tượng cúi xuống nhìn một khe nứt lớn đang lan ra dưới chân: ``Ghê thật! Nơi này thậm chí cũng không thể trụ vững được lâu nữa đâu!''

Cô quay đầu bỏ chạy ngay sau câu nói đó, tung bụi đất dưới chân như trong một cảnh phim chiến tranh. Đúng lúc cô vừa kịp hét lớn: ``\textbf{Hẹn gặp các bạn vào tuần sau!}''

Một thiên thạch khổng lồ rít lên từ bầu trời, đâm thẳng xuống ngay vị trí cô vừa đứng. Một cột sáng trắng xóa bao trùm cả khu vực. Mọi âm thanh bị hút sạch, mọi hình ảnh biến mất, chỉ còn một màu trắng mênh mông như thể cả thế giới bị xóa khỏi bản đồ.

Màn hình mờ dần\dots
%==========================================TRUYỆN PHỤ=========================================
\omk

\pov{Hikawa Kuon}

\begin{center}
  \uline{Trên xe buýt.}
\end{center}

Trước mặt tôi là một khung hình kết thúc thước phim, với hình ảnh A-san đang khoanh tay làm dấu chữ X, miệng cười khổ, với dòng chú thích được đặt ở hai bên: ``Thảm họa của sản xuất là đây.''

\img{s-22c}

Tôi chỉ biết lầm bầm: ``\dots Lại gì nữa đây?''

Về cơ bản, tôi vừa bị buộc phải xem xong một bảng xếp hạng nữa, tập hợp từ các đoạn phim cũ của \textit{hologra}, lần này mang chủ đề ``tận thế.'' Và như mọi lần, tôi biết trước kết quả sẽ ra sao.

Một bảng xếp hạng tôi thấy không chuẩn xác cho lắm.

Nói thật thì, với tôi, ngày nào ở bên họ cũng có thể được coi là một phiên bản nhỏ của ngày tận thế. Không có một buổi làm việc nào trôi qua yên ổn. Cũng không có thứ gì gọi là ``bình thường.'' Nếu tôi là dân mới vào nghề, chắc tôi đã xin nghỉ việc từ tuần thứ hai rồi. Nhưng sau hai năm rưỡi sống sót với những trò oái oăm của họ, tôi đã miễn nhiễm với khái niệm ``hỗn loạn.''

Mà nếu nói cho công bằng, cái gọi là ``tận thế'' đúng nghĩa thì tôi xin ``vinh dự'' đề cử thêm những lần có trò Old Maid. Không chỉ một mà tận hai lần\footnote{Đó là các lần ở {\flaregothic ÉPISODE 99}, quyển số Tám và ở {\flaregothic ÉPISODE 109}, quyển số Chín.}, chỉ một bộ bài và một lời nói dối nhỏ của ai đó đã làm chúng tôi phải đền gần nửa cái văn phòng, khiến cho tôi cấm túc hoàn toàn trò chơi này. Máy lạnh, máy quay, một nửa hệ thống ánh sáng, thậm chí cả máy pha cà phê cũng không sống sót nổi.

Tôi khoanh tay, gật đầu một cách mệt mỏi: ``Nhưng nếu xét theo `tận thế' theo nhiều nghĩa khác nhau thì\dots\ phải công nhận là nó cũng đúng ở một khoản nào đó. Mà, cái này là fan bình chọn mà, ắt hẳn phải có lý do gì chứ.''

Tôi thở một tiếng dài hơn mức cần thiết. Tôi không biết phải bình luận gì với cái đoạn phim vừa rồi. Mỗi lần tôi định phản đối điều gì, tôi lại tự nghe thấy câu này vang trong đầu: ``Đây là \hll\ mà.''

Một câu nói tôi đã lặp đi lặp lại còn nhiều hơn số lần tôi phải trả lời email khẩn từ hội đồng kỹ thuật, và tôi biết - không, tôi chấp nhận - là mình không thể thoát khỏi nó được nữa.

Tôi còn cách nào khác ngoài việc sống chung với lũ đâu?

\dots À mà, tôi chưa nói là cái phông nền tận thế đằng sau Sora-chan lúc dẫn chương trình chỉ là một đoạn CGI miễn phí tôi nhặt được từ một trang chia sẻ mô hình đâu nhỉ? Không rõ là ai thêm vào cái quả thiên thạch, nhưng nhìn lại thì\dots\ cũng hợp mắt đấy.
\end{document}